\section[METODOLOGÍA]{Metodología}

Este capítulo describe la metodología para desarrollar y validar el middleware de orquestación integrado con ELAN. La metodología se organiza en diez apartados: enfoque y tipo de trabajo, técnicas de recolección de información, técnicas de análisis de la información, herramientas y stack tecnológico, proceso de desarrollo (ciclo de vida completo), contratos de comunicación, despliegue y configuración del entorno, estructura de directorios del sistema, guía de replicación y uso ético de herramientas de inteligencia artificial generativa. Para evitar ambigüedades, en este capítulo se utiliza el término \textit{apartado} para referirse a la estructura metodológica general y el término \textit{fase} solo para los incrementos del proceso de desarrollo.

El apartado del proceso de desarrollo recorre el ciclo de vida completo del software: levantamiento de requerimientos, análisis y diseño, implementación en cinco fases (base del middleware, registro de modelos e instalación automática, Docker Runner e inferencia, cola de jobs y métricas, e integración con la interfaz de ELAN), pruebas y validación final, y despliegue. Esta organización permite diferenciar la estructura expositiva del capítulo respecto de los incrementos técnicos ejecutados durante el desarrollo.

El objetivo de esta metodología no es solo enumerar actividades, sino también explicar cómo el problema de integración se convirtió en una arquitectura coherente. El sistema desarrollado conecta una aplicación de escritorio en Java, un middleware en Python, servicios HTTP, contenedores Docker y backends de inferencia; por ello, cada decisión metodológica se justifica según su aporte a la interoperabilidad, la separación de responsabilidades, la validación de contratos y la reducción de dependencias entre componentes. Con el fin de mantener una lectura clara, el capítulo explica las decisiones en un nivel conceptual; el lector que requiera mayor profundidad técnica (esquemas completos, archivos de configuración, procedimientos de compilación y ejemplos íntegros de respuestas) puede dirigirse a los Anexos I, II, III y IV, que se referencian en los puntos correspondientes.

\subsection{Enfoque y tipo de trabajo}

El trabajo es de naturaleza aplicada y tecnológica. No se trata de proponer nuevas ideas en aprendizaje automático, ni de aportar un análisis lingüístico de la LSEc, sino de diseñar, implementar y validar un middleware que resuelve un problema concreto de integración: lograr que ELAN solicite inferencias a servicios de IA externos y reciba segmentos temporales compatibles con su flujo de anotación, sin alterar la arquitectura general de la herramienta y sin obligar al investigador a transcribir manualmente los resultados. Vale aclarar que el término \textit{middleware} se emplea aquí en su acepción de ingeniería de software, es decir, como una capa intermedia e independiente que comunica dos aplicaciones heterogéneas, tal como se fundamentó en el marco teórico; no se trata de un módulo interno de ELAN ni de una librería incorporada a la herramienta.

Esta denominación merece justificarse antes de describir la metodología, tomando como base la distinción entre módulo, componente y middleware presentada en el marco teórico. El software construido en este trabajo no comparte el proceso, el lenguaje ni el ciclo de vida de ELAN, por lo que no puede considerarse un módulo; tampoco se despliega ni se ejecuta junto con la herramienta, condición que caracteriza a un componente. Se trata de un proceso autónomo: se instala y se actualiza por separado, corre en su propio contenedor y atiende solicitudes a través de una interfaz de red, ubicado entre ELAN, escrito en Java, y los servicios de inferencia, escritos en Python, sin que ninguno de los dos extremos necesite conocer su implementación interna. La única pieza del trabajo que sí constituye un componente en sentido estricto es el reconocedor Java añadido a ELAN, porque vive dentro del proceso de la herramienta. Por estas razones, en el resto del documento la solución se nombra como middleware.

Como metodología de desarrollo de software se adoptó el modelo iterativo e incremental, que toma de los marcos ágiles, como Scrum, la idea de ciclos cortos con revisión y adaptación continuas, aunque sin las ceremonias y roles de equipo propios de esos marcos, dado que el TIC fue ejecutado por un solo desarrollador. El ciclo de vida cubrió todas las etapas clásicas del desarrollo: levantamiento de requerimientos, análisis y diseño, implementación, pruebas y despliegue; la diferencia frente a un modelo en cascada es que la implementación y las pruebas se repartieron en incrementos verificables: primero la API base, después el registro de modelos, la ejecución de contenedores, la cola de trabajos, la integración con ELAN y, al final, la validación de escenarios de fallo. Este enfoque resultó adecuado porque al inicio no se conocían todas las restricciones de integración entre ELAN, AVATecH, el cliente Java y el servidor del middleware; cada iteración sirvió para confirmar una decisión técnica, detectar incompatibilidades y ajustar partes ya implementadas sin rediseñar el sistema completo.

Un ejemplo de cambio iterativo es la compatibilidad HTTP entre ELAN y el middleware. Durante las pruebas de integración, el cliente Java necesitaba comunicarse de manera confiable a través de HTTP/1.1, por lo que se ajustó la configuración del servidor ASGI que ejecuta FastAPI para evitar incompatibilidades en las solicitudes. Hallazgos de este tipo muestran que las fases de desarrollo no fueron etapas cerradas, sino ciclos con retroalimentación, en los que cada validación podía obligar a revisar decisiones previas sobre comunicación, configuración o manejo de errores.

\subsection{Técnicas de recolección de información}

Para evaluar los requisitos, las restricciones y las decisiones de implementación, se utilizaron tres métodos de recopilación de información: análisis documental, ingeniería inversa del código fuente de ELAN y revisión de trabajos relacionados. Estos métodos se aplicaron en paralelo. El análisis documental permitió comprender las especificaciones y tecnologías del sistema; la ingeniería inversa reveló los puntos de extensión reales en ELAN; y la revisión de trabajos relacionados permitió analizar cómo se producen los resultados de los sistemas de reconocimiento de señas y cómo deben transformarse para ser útiles en una herramienta de anotación.

\subsubsection{Análisis documental}

La primera técnica consistió en analizar la documentación técnica. Esta revisión no se realizó solo como una actividad descriptiva, sino también como referencia para decidir qué debía implementarse en ELAN, qué debía delegarse al middleware y qué debía resolverse en los backends en contenedores.

El documento clave fue la especificación técnica de AVATecH, que describe cómo ELAN puede integrar componentes de reconocimiento externos \cite{auer2014avatech}. Esta especificación explica cómo un reconocedor obtiene información de configuración, accede a los recursos multimedia del archivo EAF, reporta progreso y devuelve fragmentos candidatos para que ELAN los agregue a los niveles de anotación. De esta revisión se concluyó que la integración debía respetar interfaces Java como \textbf{Recognizer2} o \textbf{IncrementalRecognizer}; así, las decisiones tecnológicas respondieron a las necesidades del sistema y no a preferencias de implementación.

También se revisó la documentación de las tecnologías utilizadas en la capa de orquestación. FastAPI permitió organizar los endpoints, el ciclo de inicialización y la documentación automática de la API \cite{fastapi2026features}; Pydantic habilitó modelos de datos ejecutables para validar solicitudes, respuestas y manifiestos \cite{pydantic2026models}. Docker y su SDK para Python permitieron gestionar imágenes, contenedores, redes y montajes desde el middleware \cite{docker2026overview,docker_sdk_python}; y el estilo REST orientó la definición de operaciones HTTP uniformes para consultas, instalaciones, cambios de estado y consulta de resultados \cite{fielding2000rest}. Esta revisión reforzó el vínculo entre la selección tecnológica y los requisitos del sistema.

\subsubsection{Ingeniería inversa del código fuente de ELAN}

El segundo enfoque fue la ingeniería inversa del código fuente de ELAN. La versión 7.1 se utilizó como referencia para el análisis y las pruebas de integración. La revisión no buscó rehacer ninguna parte de la herramienta, sino identificar los puntos donde ELAN permite incluir reconocedores externos en el marco AVATecH. Por esa razón, se analizó el paquete \textbf{mpi.eudico.client.annotator.recognizer}, donde se ubican las interfaces, clases de soporte y mecanismos de comunicación asociados con los reconocedores \cite{auer2010elan,auer2014avatech}.

El primer elemento identificado fue el contrato que un reconocedor externo debe cumplir. En ELAN, este contrato está representado principalmente por las interfaces \textbf{Recognizer2} e \textbf{IncrementalRecognizer}. Estas interfaces permiten recibir parámetros, iniciar el procesamiento mediante \textbf{start()}, abortar la ejecución y comunicarse con ELAN durante el proceso. La comunicación se realiza mediante \textbf{RecognizerHost}, interfaz que permite informar progreso, reportar errores y enviar segmentaciones al entorno de anotación \cite{auer2014avatech}. Esto dejó claro que el cliente Java debía trabajar dentro del flujo de ELAN y no como una herramienta externa desconectada.

El segundo elemento fue el mecanismo de entrega de parámetros. La revisión mostró que ELAN no envía una estructura única y cerrada, sino pares clave-valor a través de \textbf{setParameterValue()}; por ello se definieron parámetros explícitos para la ruta del video, el modelo seleccionado, el nivel de anotación destino y las opciones de inferencia. De esta manera, el panel de configuración no depende de valores fijos para enviar información al middleware.

El tercer elemento fue el retorno de resultados. ELAN no espera una lista genérica de predicciones, sino un objeto \textbf{Segmentation} asociado a un archivo multimedia y a un nivel de anotación, compuesto por objetos \textbf{Segment} con tiempo de inicio, tiempo de fin y etiqueta. Este formato determinó el contrato de respuesta del middleware: cada resultado del backend debía convertirse en un segmento temporal expresado en milisegundos y asociarse al nivel de anotación indicado por el usuario, pues de lo contrario ELAN no puede ubicar la anotación en la línea de tiempo.

También se revisó la construcción del panel gráfico del reconocedor. En ELAN, los paneles de configuración se basan en \textbf{AbstractRecognizerPanel}; su función es organizar controles de entrada antes de iniciar el procesamiento, no ejecutar la inferencia. Por ello se definió \textbf{AIOrchestrationRecognizerPanel} para exponer selección de modelo, validación del video, nivel de anotación destino y parámetros de inferencia. Con esta separación, el panel recopila configuración y el reconocedor ejecuta la comunicación con el middleware y transforma la respuesta en objetos compatibles con ELAN.

Finalmente, se revisó la comunicación HTTP disponible desde Java. Esta revisión orientó la implementación hacia \textbf{java.net.http.HttpClient}, porque permite construir un cliente dentro del mismo proceso de ELAN sin introducir procesos auxiliares. A partir de esta decisión, las clases de integración se ubicaron en el paquete \textbf{mpi.eudico.client.annotator.recognizer.ai}. La ubicación no responde solo a organización del código: permite reutilizar inicialización, configuración, cancelación, reporte de progreso, manejo de errores y entrega de segmentaciones mediante \textbf{RecognizerHost}, manteniendo compatibilidad con el mecanismo de extensión de AVATecH \cite{auer2014avatech}.

\subsubsection{Revisión de trabajos relacionados y modelos existentes}

La tercera técnica fue la revisión de trabajos previos sobre el reconocimiento automático de LSEc y la integración de detectores en ELAN. Esta revisión ayudó a distinguir entre el resultado de un modelo y una anotación que un investigador puede usar: un modelo puede producir una etiqueta, una predicción o una lista de probabilidades, pero ELAN necesita segmentos con inicio, fin y valor textual para agregarlos al nivel de anotación. El trabajo de Guevara sobre el intérprete de frases LSEc muestra que los sistemas de reconocimiento aprenden información visual y generan una clasificación del signo o frase observada \cite{guevara2022lsec_interprete}. Este antecedente confirmó que la salida de un reconocedor requiere una transformación adicional antes de integrarse en la línea de tiempo de ELAN.

Desde una perspectiva técnica, la revisión orientó el contrato de respuesta del middleware. Dado que los backends procesan video, la respuesta debía conservar la información temporal del archivo y no limitarse a un glosario final. Así, se definió \textbf{media\_info} para los metadatos del video y \textbf{predictions} para transportar alternativas de clasificación con sus probabilidades cuando el backend las proporciona. La revisión final queda en manos del investigador: el sistema propone segmentos y etiquetas, pero no convierte las predicciones en anotaciones definitivas sin validación humana.

La revisión también ayudó a diferenciar tres ideas que suelen confundirse: reconocimiento, segmentación y anotación. El reconocimiento produce una etiqueta candidata; la segmentación determina el intervalo de tiempo en el que ocurre el evento; y la anotación es la inserción del resultado en ELAN, siempre sujeta a revisión. Esta distinción justificó que el middleware no entregara solo una predicción textual, sino segmentos con tiempos, etiquetas y alternativas cuando existan. Además, mantiene el contrato independiente de la arquitectura de IA y de la implementación interna del backend.

Finalmente, se revisaron antecedentes de integración entre ELAN y detectores automáticos. Auer et al. describen extensiones hacia la anotación semiautomática de audio y video \cite{auer2010elan}. Este antecedente confirmó que la conexión entre ELAN y procesos externos ya existía, pero también evidenció que el presente trabajo necesitaba una capa adicional: no bastaba con invocar al detector; también había que gestionar los modelos instalados, los contenedores, los contratos HTTP, los errores y los resultados temporales. De ahí que el middleware se plantee como un orquestador entre ELAN y los backends de inferencia.

\subsection{Técnicas de análisis de la información}

La información recopilada se analizó mediante tres técnicas: descomposición funcional, diseño basado en contratos y validación incremental. Estas técnicas convirtieron la información dispersa en decisiones de diseño verificables. La descomposición determinó qué responsabilidades debía asumir cada componente; el diseño basado en contratos definió cómo se comunicarían componentes escritos en tecnologías distintas; y la validación incremental estableció cómo verificar que cada incremento siguiera funcionando dentro del flujo completo.

\subsubsection{Descomposición funcional}

La descomposición funcional consistió en dividir la integración entre ELAN, el middleware y los backends en responsabilidades pequeñas y verificables. Esta decisión se relaciona con la mantenibilidad del software, ya que un sistema con componentes cohesionados y bajo acoplamiento es más fácil de modificar, probar y extender \cite{iso25010}. En este trabajo, la descomposición se refinó durante el desarrollo; cada fase permitió identificar nuevos puntos de integración, errores de comunicación y necesidades de validación.

En este sentido, se definieron servicios internos del middleware. El \textbf{ModelRegistryService} es responsable de la instalación, validación, registro y consulta de modelos utilizando el registry.json. El \textbf{DockerLifecycleService} es responsable de la creación de imágenes y el inicio inicial durante la instalación, mientras que el \textbf{DockerService} es responsable de ejecutar los contenedores. Esta separación respondió a que instalar un modelo y ejecutar una inferencia son procesos distintos: la instalación crea los artefactos y deja el backend disponible por primera vez, mientras que la inferencia localiza el contenedor activo, envía la solicitud y gestiona el resultado. En ambos casos se emplean contenedores para aislar las aplicaciones y sus dependencias \cite{docker2026overview}.

\textbf{JobService} se definió como coordinador de inferencias: recibe la solicitud, consulta el modelo registrado, prepara la ejecución, invoca el backend y devuelve una respuesta compatible con ELAN. \textbf{JobQueue} gestiona la concurrencia mediante una cola FIFO con un límite de inferencias simultáneas, de modo que los recursos de CPU, memoria y GPU no se saturen. Por su parte, \textbf{MetricsService} mantiene contadores de observabilidad: trabajos completados, fallidos, en cola o con tiempo de espera agotado. Con esta separación es posible evaluar el estado del sistema sin mezclar métricas, ejecución y persistencia.

\subsubsection{Diseño dirigido por contratos}

La segunda técnica se basó en contratos. En este trabajo, un contrato es la definición explícita de los datos que un componente acepta, valida y produce antes de comunicarse con otro componente. Esta decisión fue necesaria porque ELAN está implementado en Java, el middleware en Python y los backends se exponen como servicios HTTP. En una arquitectura de este tipo, asumir estructuras implícitas aumenta el riesgo de errores que son difíciles de diagnosticar. Por lo tanto, se definieron contratos para solicitudes HTTP, respuestas de inferencia y paquetes de instalación usando \textbf{manifest.json}, según el principio REST de interfaces uniformes y mensajes bien definidos \cite{fielding2000rest}.

En el middleware, los contratos se implementaron utilizando modelos Pydantic junto con FastAPI. Esta combinación permitió la validación de cuerpos de solicitud antes de ejecutar la lógica de negocio. Si el cliente Java envía un campo faltante, un tipo incorrecto o una estructura inválida, entonces el middleware puede enviar un error estructurado en lugar de pasar por la falla interna. Por tanto, los contratos no quedaron solo como documentación: se convirtieron en una parte ejecutable del sistema \cite{pydantic2026models,fastapi2026features}.

El diseño por contratos también se aplicó al paquete de cada modelo. El archivo \textbf{manifest.json} describe identificador, versión, artefactos, forma de construcción del backend Docker y endpoints de inferencia. Validarlo de forma temprana permite detectar paquetes incompletos antes de construir o iniciar un contenedor, lo que resulta preferible a descubrir el problema durante una inferencia, cuando el modelo ya figura como registrado y el usuario espera que esté disponible.

\subsubsection{Validación incremental mediante pruebas funcionales}

La tercera técnica fue la validación mediante pruebas funcionales. Cada fase produjo un incremento verificable: API base, registro de modelos, gestión de contenedores, cola de inferencia, integración con ELAN y recuperación ante fallos. Se eligió este tipo de validación porque el sistema puede fallar por causas muy distintas entre sí: errores de contrato HTTP, rutas de video inválidas, contenedores detenidos, tiempos de espera del backend o incompatibilidades entre la respuesta del middleware y el formato esperado por ELAN.

Las pruebas se realizaron sobre las interfaces públicas. En el middleware se enviaron solicitudes HTTP válidas e inválidas y las respuestas se compararon con los contratos definidos; en ELAN se comprobó que los segmentos generados aparecieran en el nivel de anotación configurado, con tiempos de inicio y fin consistentes con la respuesta del backend. De este modo se validó el flujo completo: cliente Java, API REST, servicios internos, backend en contenedor y retorno de segmentaciones a ELAN. Además, cualquier escenario puede repetirse siempre que se conserven la solicitud, el paquete y las condiciones de ejecución.

En conjunto, las tres técnicas llevaron de información técnica dispersa a una arquitectura verificable: la descomposición funcional organizó las responsabilidades, el diseño por contratos definió interfaces comprobables y la validación incremental confirmó que esas decisiones funcionaban en el sistema completo.

\subsection{Herramientas y stack tecnológico}

La Tabla~\ref{tab:stack_tecnologico} resume las herramientas utilizadas y las versiones establecidas como base durante el desarrollo. Estas versiones no son la única combinación posible, pero sí el entorno probado para evitar incompatibilidades entre el cliente Java de ELAN, el middleware, Docker y el backend de inferencia. La selección tecnológica tuvo en cuenta los siguientes puntos: comunicación explícita entre componentes, validación de contratos y aislamiento de dependencias a través de contenedores.

\begin{table}[H]
\centering
\caption{Stack tecnológico del sistema desarrollado}
\label{tab:stack_tecnologico}
\footnotesize
\renewcommand{\arraystretch}{1.18}
\setlength{\tabcolsep}{5pt}

\begin{tabular}{
    >{\raggedright\arraybackslash}p{3.1cm}
    >{\raggedright\arraybackslash}p{3.4cm}
    >{\raggedright\arraybackslash}p{6.2cm}
}
\toprule
\textbf{Componente} & \textbf{Tecnología y versión} & \textbf{Razón principal de selección} \\
\midrule

Middleware &
Python 3.11 &
Compatibilidad con bibliotecas de inteligencia artificial y mejoras de rendimiento del intérprete. \\

Framework API &
FastAPI 0.115 &
Definición de endpoints REST, validación integrada y documentación OpenAPI. \\

Validación de datos &
Pydantic 2.x &
Modelos de datos ejecutables para validar solicitudes, respuestas y manifiestos. \\

Servidor ASGI &
Uvicorn 0.32 &
Ejecución de la aplicación FastAPI sobre el estándar ASGI. \\

Contenerización &
Docker 25.x y Docker Compose &
Aislamiento de modelos, dependencias y redes de ejecución. \\

SDK Docker &
Docker SDK for Python 7.x &
Control programático de imágenes y contenedores desde el middleware. \\

Persistencia local &
JSON sobre bind mounts &
Registro simple y auditable de modelos en un entorno local de usuario único. \\

Cliente ELAN &
Java 11 &
Disponibilidad del paquete \texttt{java.net.http} para comunicación HTTP desde ELAN. \\

Cliente--middleware &
HTTP/1.1 + JSON &
Interfaz simple entre el cliente Java y el middleware REST. \\

Middleware--backend &
HTTP interno en red Docker &
Comunicación aislada entre contenedores sin exponer todos los servicios al anfitrión. \\

\bottomrule
\end{tabular}
\end{table}

Python 3.11 fue seleccionado para implementar el middleware porque permite construir servicios de integración compatibles con el ecosistema utilizado por muchos backends de inferencia. En este proyecto, Python no se usó para entrenar modelos, sino para implementar la API, validar datos, coordinar jobs, invocar servicios externos e interpretar respuestas. La versión 3.11 se fijó como línea base por sus mejoras de rendimiento y por ofrecer una plataforma estable para bibliotecas modernas sin depender de versiones experimentales \cite{python311docs}.

FastAPI fue elegido para implementar la API REST porque permite declarar endpoints HTTP usando tipos de Python y modelos Pydantic. Esto fue importante porque el cliente Java envía solicitudes JSON y el middleware debe rechazar de forma controlada cualquier estructura incompleta o inválida. Además, FastAPI genera documentación OpenAPI y soporta carga de archivos mediante \textbf{UploadFile}, lo que facilitó la instalación de paquetes ZIP de modelos \cite{fastapi2026features}.

Se eligió Pydantic 2.x para representar solicitudes, respuestas y manifiestos como modelos de datos. Con esta decisión, errores de tipo, campos faltantes o estructuras inválidas se detectan antes de ejecutar la lógica de negocio \cite{pydantic2026models}. Por ello, Pydantic se aplicó tanto a endpoints del middleware como al archivo \textbf{manifest.json}, evitando que errores de empaquetado se descubran recién durante una inferencia.

Docker y Docker Compose se seleccionaron para aislar los backends de inferencia junto con sus dependencias, ya que mezclarlas con ELAN o con el entorno base del middleware incrementaría el riesgo de incompatibilidades. Docker empaqueta servicios en contenedores y los conecta mediante redes administradas \cite{docker2026overview}. Sobre esa base, el Docker SDK for Python se incorporó para construir imágenes, crear contenedores, iniciar servicios, inspeccionar estados y detener backends desde el propio middleware, sin depender de comandos manuales de consola \cite{docker_sdk_python}.

Se eligió Java 11 para el cliente integrado en ELAN porque incorpora \textbf{java.net.http}, paquete que proporciona \textbf{HttpClient}, \textbf{HttpRequest} y \textbf{HttpResponse} como API estándar para comunicación HTTP \cite{java11_httpclient}. Esta decisión evitó agregar librerías externas al módulo de integración. Además, \textbf{HttpClient} permitió configurar HTTP/1.1 para mantener compatibilidad con el servidor del middleware. Para instalar modelos desde ELAN se construyó manualmente el cuerpo \textbf{multipart/form-data}, manteniendo un cliente Java ligero y controlado.

La persistencia local se resolvió mediante archivos JSON almacenados en bind mounts. Esta decisión responde a un escenario local de usuario único, donde un investigador ejecuta ELAN, el middleware y los backends en su estación de trabajo. En ese contexto, una base de datos relacional exigiría una configuración innecesaria solo para registrar modelos, verificar su estado y recuperar manifiestos. El archivo \textbf{registry.json} funciona como un registro auditable y, gracias a los bind mounts, los datos permanecen en el sistema de archivos del anfitrión aunque los contenedores se recreen \cite{docker2026overview}.

Uvicorn se eligió como servidor ASGI para ejecutar la aplicación FastAPI. Su papel no es realizar inferencias, sino mantener accesible la API del middleware: mientras una inferencia prolongada ocupa un hilo de trabajo, los endpoints de salud, gestión y métricas siguen respondiendo. Este comportamiento complementa a \textbf{JobQueue}, que limita las inferencias simultáneas sin bloquear el resto del servicio. La versión 0.32 se fijó por compatibilidad con FastAPI y con la configuración utilizada durante el desarrollo \cite{uvicorn_docs}.
\vspace{1\baselineskip}

\subsection{Proceso de desarrollo: ciclo de vida del software}

El desarrollo siguió el ciclo de vida completo del software bajo un modelo iterativo e incremental. La Figura~\ref{fig:ciclo_desarrollo} resume ese ciclo: el levantamiento de requerimientos definió qué debía hacer el sistema; el análisis y diseño convirtió esos requerimientos en una arquitectura; la implementación se dividió en cinco fases que entregaron incrementos operativos; la etapa de pruebas validó el sistema completo (fase 6); y el despliegue dejó la solución instalada y documentada. Las flechas de retorno indican que el proceso no fue lineal: cada validación podía obligar a revisar el diseño o ajustar incrementos anteriores, como ocurrió con la compatibilidad HTTP entre el cliente Java y el servidor del middleware.

\begin{figure}[H]
\centering
\resizebox{0.98\textwidth}{!}{%
\begin{tikzpicture}[
    etapa/.style={rectangle, rounded corners=2pt, draw=black!70, fill=gray!12,
                  text width=2.5cm, minimum height=1.5cm, align=center, font=\footnotesize},
    flecha/.style={-{Stealth[length=2.6mm]}, thick},
    retro/.style={-{Stealth[length=2.2mm]}, thick, dashed, black!60},
    node distance=0.55cm
]
\node[etapa] (req)  {Levantamiento de requerimientos};
\node[etapa, right=of req]  (dis)  {Análisis y diseño de la arquitectura};
\node[etapa, right=of dis]  (imp)  {Implementación incremental\\ (fases 1 a 5)};
\node[etapa, right=of imp]  (val)  {Pruebas y validación final\\ (fase 6)};
\node[etapa, right=of val]  (desp) {Despliegue y documentación};

\draw[flecha] (req) -- (dis);
\draw[flecha] (dis) -- (imp);
\draw[flecha] (imp) -- (val);
\draw[flecha] (val) -- (desp);

\draw[retro] (imp.north) -- ++(0,0.6) -| (dis.north)
      node[pos=0.22, above, font=\scriptsize] {ajustes de diseño};
\draw[retro] (val.south) -- ++(0,-0.6) -| (imp.south)
      node[pos=0.22, below, font=\scriptsize] {retroalimentación de cada iteración};
\end{tikzpicture}%
}
\caption{Ciclo de vida de desarrollo seguido en el trabajo}
\label{fig:ciclo_desarrollo}
\end{figure}

La lógica del avance fue incremental: primero se consiguió que ELAN, el middleware y los backends de inferencia se comunicaran; luego se añadieron los componentes que hacen esa comunicación estable, observable y recuperable ante fallos. El criterio de progreso no era que el código compilara, sino que cada parte cumpliera su porción del contrato dentro de la arquitectura general.

\subsubsection{Levantamiento de requerimientos}

Los requerimientos del sistema se derivaron de las técnicas de recolección de información descritas antes: la especificación de AVATecH y la ingeniería inversa de ELAN definieron las condiciones que la integración debía respetar, mientras que la revisión de trabajos relacionados aclaró qué forma debían tener los resultados para ser útiles al investigador. A partir de ese material se establecieron los requerimientos funcionales y no funcionales de la Tabla~\ref{tab:requerimientos}, que guiaron el diseño de la arquitectura y sirvieron después como criterios de verificación en la validación final.

\begin{table}[H]
\centering
\caption{Requerimientos funcionales y no funcionales del sistema}
\label{tab:requerimientos}
\footnotesize
\renewcommand{\arraystretch}{1.18}
\setlength{\tabcolsep}{5pt}

\begin{tabular}{
    >{\centering\arraybackslash}p{1.4cm}
    >{\raggedright\arraybackslash}p{11.6cm}
}
\toprule
\textbf{Código} & \textbf{Descripción} \\
\midrule
RF-01 & Instalar modelos de IA desde un paquete ZIP que incluya su manifiesto de instalación. \\
RF-02 & Listar los modelos instalados y consultar el detalle de cada uno. \\
RF-03 & Activar y desactivar modelos previamente instalados. \\
RF-04 & Ejecutar inferencias sobre un video y devolver segmentos temporales con etiquetas de glosa. \\
RF-05 & Consultar el resultado de un trabajo de inferencia y las métricas acumuladas del servicio. \\
RF-06 & Gestionar el ciclo de vida de los backends en contenedores: construcción, arranque, verificación de disponibilidad y reinicio. \\
RF-07 & Permitir la configuración y ejecución del reconocedor desde la interfaz gráfica de ELAN, incluida la gestión de modelos. \\
RF-08 & Informar los errores con códigos y mensajes comprensibles para el usuario y para otros clientes HTTP. \\
\midrule
RNF-01 & Ejecución completamente local, sin enviar videos ni anotaciones a servicios externos. \\
RNF-02 & Aislamiento de las dependencias de cada modelo respecto de ELAN y del propio middleware. \\
RNF-03 & Control configurable de la cantidad de inferencias simultáneas para no saturar el equipo. \\
RNF-04 & Persistencia del registro de modelos y capacidad de recuperarlo tras reinicios o pérdidas parciales. \\
RNF-05 & Compatibilidad con ELAN 7.1 y su mecanismo de extensión AVATecH, sin alterar la arquitectura de la herramienta. \\
\bottomrule
\end{tabular}
\end{table}

\subsubsection{Análisis y diseño}

Antes de implementar el sistema se revisaron la especificación de AVATecH, el código fuente de ELAN y las restricciones operativas de los backends de IA. Esta revisión ayudó a delimitar qué parte del flujo debía quedar dentro de ELAN y qué parte correspondía a un servicio externo. AVATecH establece que los reconocedores externos deben implementarse a través de interfaces esperadas por ELAN y devolver resultados convertibles en anotaciones temporales \cite{auer2014avatech}; por lo tanto, la solución debía cubrir la configuración, el reporte de progreso, la cancelación y el retorno de segmentos.

Se analizaron cuatro alternativas de integración antes de definir la arquitectura. La primera fue ejecutar procesos externos desde ELAN mediante línea de comandos o archivos intermedios. Aunque mantenía el modelo fuera de ELAN, complicaba el reporte de progreso, la cancelación y la conversión de resultados a estructuras compatibles con el anotador. La segunda fue empaquetar toda la lógica como una librería JAR incorporada a ELAN, de manera que la herramienta resolviera la inferencia dentro de su propio proceso. Esta opción se descartó por una razón de fondo: los modelos de reconocimiento no son código Java, sino pipelines en Python con dependencias nativas de visión artificial y aprendizaje profundo, pesos de cientos de megabytes y requisitos de hardware propios; nada de eso puede empaquetarse en un JAR. Además, una librería comparte el ciclo de vida de la aplicación que la carga: cada actualización de un modelo obligaría a recompilar y redistribuir ELAN, y un fallo o fuga de memoria del modelo comprometería directamente la estabilidad de la herramienta de anotación. La tercera alternativa fue tender un puente entre Java y las bibliotecas nativas mediante JNI, pero introducía una frontera compleja entre Java, Python, bibliotecas nativas y dependencias del modelo \cite{oracle_jni}, y reducía la portabilidad, porque el usuario tendría que mantener un entorno compatible con todas esas dependencias.

La alternativa adoptada fue implementar en ELAN un reconocedor Java que actúe como cliente HTTP de un servicio local. Con ello el sistema queda organizado como cliente-servidor: el middleware es un software independiente que se instala, arranca y actualiza por separado, y ELAN se limita a consumir su API. La elección de HTTP merece una justificación explícita porque el escenario previsto es local y podría parecer innecesario usar un protocolo de red. En la práctica, el costo es despreciable (la comunicación viaja por la interfaz de loopback de la misma máquina, sin salir a la red) y las ventajas son concretas: los procesos quedan aislados, de modo que un fallo del modelo no afecta a ELAN; la frontera entre Java y Python es explícita y verificable, lo que permitió probar el middleware con la herramienta Postman sin depender de ELAN; los modelos se actualizan sin recompilar la herramienta; y la misma interfaz sirve sin cambios si en el futuro el servidor se ubica en otra máquina de la red. La comunicación HTTP/1.1 con JSON se eligió por ser explícita, verificable y desacoplada, de acuerdo con el estilo REST \cite{fielding2000rest}. Para el cliente Java se usó \textbf{java.net.http.HttpClient}, disponible desde Java 11, evitando añadir librerías externas al código de ELAN \cite{java11_httpclient}.

A partir de esta decisión se definió una arquitectura de tres capas: ELAN como cliente de anotación, el middleware como orquestador local y los backends como servicios contenerizados. Docker permitió empaquetar cada backend con sus dependencias y ejecutarlo de forma aislada del sistema anfitrión \cite{docker2026overview}. Así, el investigador puede seguir trabajando en ELAN, mientras el middleware instala modelos, verifica disponibilidad, crea contenedores y envía solicitudes de inferencia.

\begin{figure}[htbp]
    \centering
    \includegraphics[width=0.85\textwidth]{Diagramas/diagrama_arquitectura_general.png}
    \caption{Arquitectura general del sistema de orquestación}
    \label{fig:arquitectura_general}
\end{figure}

La Figura~\ref{fig:arquitectura_general} resume esta separación de responsabilidades. ELAN no ejecuta directamente el pipeline de IA; envía una solicitud al middleware. El middleware valida la petición, localiza el modelo, coordina Docker y devuelve segmentos temporales que ELAN puede insertar en el nivel de anotación configurado.

En esta etapa también se definió el contrato de instalación mediante \textbf{manifest.json}. Este archivo acompaña a cada paquete de modelo y declara identidad, versión, tarea, artefactos, runtime, configuración del contenedor, endpoints del backend y opciones visibles en ELAN. Definir el esquema antes de implementar la instalación permitió validar los paquetes de forma temprana y devolver errores estructurados, aprovechando las capacidades de Pydantic y FastAPI \cite{pydantic2026models,fastapi2026features}.

El diseño también estableció la estructura interna de clases del middleware, que se mantuvo estable durante todas las fases de implementación. La Figura~\ref{fig:diagrama_clases} presenta el diagrama de clases de los servicios principales con sus operaciones públicas y sus relaciones de uso. La capa de API recibe las solicitudes y las delega en dos coordinadores: \textbf{ModelRegistryService}, a cargo de la instalación y el registro de modelos, y \textbf{JobService}, a cargo de las inferencias. Ambos se apoyan en servicios especializados: la cola de trabajos, las métricas, la selección del mecanismo de ejecución y los servicios que hablan con Docker. La clase abstracta \textbf{BaseRunner} define el contrato de ejecución y \textbf{DockerRunner} lo implementa para backends en contenedores, de modo que otros mecanismos de ejecución podrían agregarse sin tocar el resto del sistema. Las clases Java integradas en ELAN se describen en el apartado de estructura de directorios y su inventario completo consta en el Anexo~IV.

% El diagrama de clases se elabora en draw.io y se exporta como PNG
% (archivo fuente: Diagramas/diagrama_clases_middleware.drawio).
\begin{figure}[H]
    \centering
    \includegraphics[width=0.97\textwidth]{Diagramas/DiagramaClasesMiddleware.drawio.png}
    \caption{Diagrama de clases de los servicios principales del middleware}
    \label{fig:diagrama_clases}
\end{figure}

\subsubsection{Fase 1 — Base del middleware}

El primer paso fue construir una base del middleware simple, ejecutable y verificable. Antes de abordar la gestión de modelos o la inferencia, se implementó un servicio HTTP local capaz de iniciarse en Docker, cargar su configuración, exponer un endpoint de salud y generar registros de diagnóstico. Comenzar por esta base redujo el riesgo inicial, pues comprobó que Python, FastAPI, Uvicorn, Docker y la red de despliegue funcionaban en conjunto.

Desde el principio, las responsabilidades internas se separaron. La capa \textbf{api/} contiene los routers HTTP; \textbf{services/} concentra la lógica de negocio; \textbf{runners/} agrupa los mecanismos de ejecución; \textbf{schemas/} define los contratos Pydantic; y \textbf{core/} centraliza configuración y logging. Esta organización evita que los handlers acumulen lógica de orquestación y facilita la modificación de una parte del sistema sin afectar a todo el middleware.

La configuración se centralizó en \textbf{app/core/config.py} mediante una clase \textbf{Settings} basada en Pydantic. Las rutas, los límites de concurrencia, los tiempos de espera, el nivel de registro y los parámetros del contenedor se declaran allí. El registro estructurado se implementó en \textbf{app/core/logging\_config.py}, con mensajes que incluyen severidad, marca de tiempo, módulo y descripción. La verbosidad se controla con \textbf{MIDDLEWARE\_LOG\_LEVEL}, siguiendo la práctica de configurar el comportamiento desde el entorno \cite{python_logging}.

El manejo global de errores se definió en \textbf{main.py} a través de manejadores de excepciones de FastAPI para errores de registro, Docker y selección de runners. Todos los errores se convierten en una respuesta uniforme con código semántico y detalle legible.

\begin{center}
\begin{minipage}{0.85\textwidth}
\hrule
\vspace{0.12cm}
\small
\begin{Verbatim}[baselinestretch=0.82]
{"error_code": "CODIGO_DESCRIPTIVO", "detail": "Mensaje legible."}
\end{Verbatim}
\vspace{0.05cm}
\hrule

\captionof{codigo}{Formato uniforme de respuesta de error del middleware}
\label{lst:error_response}
\end{minipage}
\end{center}

Este formato fue importante para la integración con ELAN, porque el cliente Java puede leer siempre el campo \textbf{detail}, sin importar si el fallo proviene del registro de modelos, Docker o el backend de inferencia.

La configuración del contenedor del middleware se definió con un \textbf{Dockerfile} basado en \textbf{python:3.11-slim} y un \textbf{docker-compose.yml} con servicio, bind mounts, variables de entorno y red compartida. La red \textbf{elan-ai-shared} se declaró con nombre fijo para que los contenedores de modelos instalados posteriormente se unan al mismo espacio de comunicación \cite{docker2026overview}.

Durante el arranque, antes de aceptar solicitudes, el middleware configura los logs y escanea \textbf{data/bootstrap\_manifests/} para recuperar modelos instalados en sesiones previas. El criterio de verificación fue levantar el servicio con \textbf{docker compose up --build -d} y confirmar la respuesta de \textbf{GET \symbol{47}health}:

\begin{center}
\begin{minipage}{0.95\textwidth}
\hrule
\vspace{0.12cm}
\small
\begin{Verbatim}[baselinestretch=0.82]
{"status": "ok", "service": "elan-ai-orchestrator", "version": "0.1.0"}
\end{Verbatim}
\vspace{0.05cm}
\hrule

\captionof{codigo}{Respuesta esperada del endpoint de salud del middleware}
\label{lst:health_response}
\end{minipage}
\end{center}

Además de la respuesta HTTP, se verificó que Uvicorn iniciara sin errores y que el endpoint de salud siguiera respondiendo mientras una inferencia larga estaba en ejecución. Con ello se confirmó que la base del middleware era suficiente para el caso de uso local previsto \cite{uvicorn_docs}.

\subsubsection{Fase 2 — Registro de modelos e instalación automática}

La segunda fase abordó la gestión de modelos. El middleware debía recibir un paquete ZIP, validar su contenido, construir la imagen Docker del backend, crear el contenedor, verificar su disponibilidad y registrar el modelo instalado. Dado que esta fase combina sistema de archivos, validación, Docker y persistencia, se diseñó como un flujo secuencial con puntos de control y recuperación.

El proceso comienza con la recepción del ZIP y la localización del \textbf{manifest.json}. El middleware valida el manifiesto contra el esquema \textbf{ModelManifest}, comprueba que los artefactos declarados existan dentro del paquete, rechaza rutas inseguras y verifica que la combinación de \textbf{model\_id} y \textbf{version} no esté ya instalada. Solo entonces extrae el paquete, construye la imagen Docker, crea el contenedor y ejecuta la comprobación de disponibilidad. Si la construcción, el arranque, esa comprobación o la persistencia fallan, se ejecuta un rollback que elimina los archivos extraídos y la imagen construida, de modo que el sistema regresa a su estado anterior.

La validación previa a la extracción actúa como barrera de seguridad y consistencia: el identificador del modelo solo admite caracteres seguros, todos los artefactos declarados deben existir dentro del ZIP, las rutas no pueden contener \textbf{../} ni ser absolutas, y la tarea debe estar entre las soportadas. Estas reglas evitan instalar paquetes incompletos o capaces de escribir fuera del directorio previsto, y se apoyan en la validación temprana descrita en el diseño por contratos.

El \textbf{manifest.json} es el contrato entre quien empaqueta un modelo y el middleware. La Tabla~\ref{tab:manifest_campos} resume los campos que el orquestador necesita para instalar, ejecutar y presentar un modelo sin codificar esos datos dentro del middleware. El esquema completo del archivo, con los valores admitidos en cada objeto y el manifest real del modelo de referencia, se recoge en el Anexo~I.

\begin{table}[H]
\centering
\caption{Campos principales del \texttt{manifest.json}}
\label{tab:manifest_campos}
\small
\renewcommand{\arraystretch}{1.15}
\begin{tabular}{@{}>{\raggedright\arraybackslash}p{3.1cm}>{\raggedright\arraybackslash}p{2.1cm}>{\raggedright\arraybackslash}p{7.2cm}@{}}
\toprule
\textbf{Campo} & \textbf{Tipo} & \textbf{Descripción} \\
\midrule
\texttt{model\_id} & string & Identificador único del modelo. Solo admite caracteres \texttt{[A-Za-z0-9\_.-]}. \\
\texttt{version} & string & Versión semántica del modelo, por ejemplo \texttt{1.0.0}. \\
\texttt{task} & string & Tarea del modelo: \texttt{video\_segmentation} o \newline\texttt{video\_segmentation\_and}\newline\texttt{\_gloss\_classification}. \\
\texttt{runtime.mode} & string & Modo de ejecución. En esta implementación, el único valor soportado es \texttt{docker}. \\
\texttt{artifacts} & object & Rutas relativas dentro del ZIP para cada artefacto declarado; todos los archivos deben existir. \\
\texttt{backend\_config} & object & Nombre de imagen, puerto, ruta de comprobación de disponibilidad y timeout de arranque usados durante la instalación. \\
\texttt{container} & object & Puerto, ruta de comprobación de disponibilidad y ruta de inferencia usados durante cada ejecución. \\
\texttt{ui} & object & Configuración de la interfaz de ELAN: nivel de anotación destino, etiqueta y modo de etiquetado. \\
\bottomrule
\end{tabular}
\end{table}

La estructura esperada del ZIP forma parte del mismo contrato. Cualquier backend nuevo debe respetar esta organización para que el middleware ubique el Dockerfile, las dependencias, los pesos, el vocabulario y la configuración del pipeline sin reglas especiales por modelo:

\begin{center}
\begin{minipage}{0.88\textwidth}
\hrule
\vspace{0.12cm}
\small
\begin{Verbatim}[baselinestretch=0.82]
mi_modelo.zip
|-- manifest.json              (obligatorio, en la raiz)
|-- backend/
|   |-- Dockerfile             (imagen Docker del backend)
|   |-- requirements.txt
|   `-- app/
|       |-- main.py            (endpoints /health e /infer)
|       `-- ...
|-- weights/                   (pesos del modelo)
|-- vocab/                     (vocabulario de glosas)
`-- config/
    `-- pipeline_config.json   (configuracion del pipeline)
\end{Verbatim}
\vspace{0.05cm}
\hrule

\captionof{codigo}{Estructura esperada del paquete ZIP de un modelo}
\label{lst:estructura_zip_modelo}
\end{minipage}
\end{center}

Para mejorar la resiliencia se implementaron \textit{bootstrap manifests}. Durante la instalación, el middleware guarda una copia del manifest en \textbf{data/bootstrap\_manifests/} con el nombre \textbf{<model\_id>\_\_<version>.json}. Al iniciar, el middleware escanea ese directorio y vuelve a registrar los modelos si \textbf{registry.json} fue eliminado, siempre que el directorio persistente del anfitrión siga disponible.

El nombre del contenedor Docker se genera de forma determinística a partir del identificador y la versión del modelo:

\begin{center}
\begin{minipage}{0.55\textwidth}
\hrule
\vspace{0.12cm}
\small
\begin{Verbatim}[baselinestretch=0.82]
elan-ai-model-{model_id}-{version}
\end{Verbatim}
\vspace{0.05cm}
\hrule

\captionof{codigo}{Patrón de nombre del contenedor Docker de un modelo}
\label{lst:nombre_contenedor_modelo}
\end{minipage}
\end{center}

Con \textbf{model\_id=lsec\_bio\_gloss\_final\_v1} y \textbf{version=1.0.0}, el nombre generado es \textbf{elan-ai-model-lsec\_bio\_gloss\_final\_v1-1.0.0}. Esta regla evita almacenar un identificador adicional: el middleware puede localizar el contenedor derivando su nombre desde los metadatos. Las operaciones de construcción, creación, inicio e inspección se realizaron con el SDK de Docker para Python \cite{docker_sdk_python}.

La Figura~\ref{fig:flujo_instalacion} muestra el flujo de instalación completo. La ruta principal termina con el modelo registrado; la rama de error representa el rollback cuando falla la construcción, el arranque, la comprobación de disponibilidad o la persistencia.

\begin{figure}[htbp]
    \centering
    \includegraphics[width=1\textwidth]{Diagramas/Proceso de instalacion de modelo.png}
    \caption{Proceso de instalación de un modelo}
    \label{fig:flujo_instalacion}
\end{figure}

También se resolvió un caso frecuente de empaquetado en Windows. Si el ZIP conserva una carpeta raíz, el manifest puede quedar en \textbf{lsec\_bio\_gloss\_final\_v1/manifest.json} en lugar de la raíz del comprimido. El middleware detecta si todas las entradas comparten el mismo prefijo y lo elimina al leer el manifest y extraer archivos. De esta manera, un detalle de empaquetado no se convierte en un error difícil de interpretar para el usuario.

El criterio de verificación fue instalar \textbf{lsec\_bio\_gloss\_final\_v1} versión \textbf{1.0.0} mediante \textbf{POST \symbol{47}api\symbol{47}v1\symbol{47}models\symbol{47}install}. Se comprobó que el modelo apareciera como \textbf{available}, que \textbf{registry.json} almacenara la entrada completa, que existiera el bootstrap manifest y que una instalación duplicada respondiera HTTP 409 con \textbf{MODEL\_ALREADY\_EXISTS}.

\subsubsection{Fase 3 — Docker Runner e inferencia}

La tercera fase introdujo \textbf{DockerRunner}, el componente que ejecuta inferencias en el contenedor del modelo. Su trabajo es preparar el entorno antes de cada solicitud: verificar que el contenedor exista, iniciarlo si está detenido, comprobar su disponibilidad, construir la carga útil, enviar la solicitud al backend y adaptar la respuesta al contrato del middleware. Con esta separación, \textbf{JobService} queda libre de los detalles del SDK de Docker.

La comunicación entre el middleware y el backend ocurre dentro de la red \textbf{elan-ai-shared}. El middleware usa el nombre del contenedor como nombre de host y deja que el DNS interno de Docker lo resuelva. Gracias a ello no es necesario exponer un puerto del anfitrión por cada modelo, se elimina la dependencia de direcciones IP y varios backends pueden convivir en la misma red \cite{docker2026overview}.

El método \textbf{ensure\_container()} controla el ciclo de vida antes de inferir. Si el contenedor ya está ejecutándose, se reutiliza; si existe pero está detenido, se reinicia; si no existe pero la imagen sí, se crea un contenedor nuevo con red, nombre y volúmenes. Si falta la imagen, se lanza \textbf{DockerImageNotFoundError}, que el manejador global transforma en HTTP 404 para indicar que el modelo debe reinstalarse. Con esto, el sistema puede recuperarse de reinicios de Docker Desktop o del equipo anfitrión.

Un punto crítico fue traducir la ruta del video. ELAN selecciona archivos del sistema anfitrión, por ejemplo \textbf{C:/Users/imbaq/Videos/seña.mp4}; el backend, al ejecutarse en Linux dentro del contenedor, accede al archivo mediante \textbf{\symbol{47}data\symbol{47}videos}. Por ello, \textbf{DockerRunner} toma el nombre del archivo recibido y lo antepone a la ruta interna del contenedor, produciendo \textbf{\symbol{47}data\symbol{47}videos\symbol{47}seña.mp4}. La variable \textbf{MIDDLEWARE\_VIDEOS\_DIR} se pasa al SDK de Docker como cadena cruda para conservar correctamente rutas de Windows \cite{docker_sdk_python}.

Antes de llamar al backend, \textbf{DockerRunner} ejecuta polling contra \textbf{GET \symbol{47}health} cada 250 milisegundos. Se diferenciaron dos timeouts: \textbf{backend\_config.startup\_timeout\_sec = 180} durante instalación, porque el backend puede cargar pesos por primera vez; y \textbf{container.startup\_timeout\_sec = 30} durante inferencia, porque el contenedor ya debería estar preparado o reiniciarse con rapidez.

Los estados de un job se resumen en la Figura~\ref{fig:estados_job}. La figura funciona como contrato operativo para interpretar logs, métricas y respuestas: todo job inicia en \textbf{RECEIVED}, pasa por validación, preprocesamiento, cola y ejecución, y termina en \textbf{COMPLETED}, \textbf{FAILED} o \textbf{TIMEOUT}.

\begin{figure}[H]
    \centering
    \includegraphics[width=0.55\textwidth]{Diagramas/Maquina de estados del Job.png}
    \caption{Máquina de estados de un job de inferencia}
    \label{fig:estados_job}
\end{figure}

Cada estado cumple una función específica: \textbf{VALIDATING} verifica el modelo; \textbf{PREPROCESSING} construye \textbf{InferenceInput} y traduce la ruta del video; \textbf{QUEUED} indica espera por el límite de concurrencia; \textbf{RUNNING} corresponde a la llamada de \textbf{DockerRunner}; \textbf{POSTPROCESSING} valida y adapta la respuesta; y los estados finales clasifican el resultado de la ejecución.

La respuesta exitosa incluye los segmentos y un campo de trazabilidad que registra el runner, la imagen Docker, el contenedor, los tiempos de validación, cola, inferencia y postprocesamiento, además del historial de estados. Esta trazabilidad facilita reconstruir la ejecución durante las pruebas y diagnosticar fallos en una arquitectura compuesta por varios servicios; el Anexo~II describe campo por campo esta información.

La Figura~\ref{fig:flujo_inferencia} relaciona los componentes de la inferencia: ELAN envía la solicitud, el middleware encola el job, DockerRunner prepara el contenedor, el backend procesa el video y el middleware adapta los segmentos para ELAN. El bloque interno del backend resume la canalización definida para el modelo utilizado en la validación; esta canalización pertenece al backend ejecutado en el contenedor y no al núcleo del middleware.

\begin{figure}[H]
    \centering
    \includegraphics[width=1\textwidth]{Diagramas/ELAN Video Segmentation-2026-06-20-052314.png}
    \caption{Flujo completo de una inferencia desde ELAN hasta las anotaciones}
    \label{fig:flujo_inferencia}
\end{figure}

Los errores de \textbf{DockerRunner} se tradujeron a códigos HTTP con significado propio: \textbf{DockerNotAvailableError} responde 503; \textbf{DockerImageNotFoundError}, 404; \textbf{DockerContainerStartError} y \textbf{DockerContainerHealthcheckFailedError}, 502; \textbf{DockerInferenceError}, 502 con el detalle reportado por el backend; y \textbf{DockerTimeoutError}, 504. De este modo, ELAN, Postman o cualquier otro cliente identifica la causa de una falla sin necesidad de revisar los logs.

El criterio de verificación fue ejecutar una inferencia desde Postman y recibir al menos un segmento con \textbf{start\_ms}, \textbf{end\_ms}, \textbf{label} y \textbf{predictions}. También se comprobó que los logs mostraran la transición de \textbf{RECEIVED} a \textbf{COMPLETED} y que \textbf{trace} incluyera tiempos por etapa e identificador del contenedor.

\subsubsection{Fase 4 — Cola de jobs, métricas y estado TIMEOUT}

La cuarta fase añadió control de concurrencia, métricas y separación explícita entre \textbf{FAILED} y \textbf{TIMEOUT}. Estos componentes no modifican el pipeline de IA; regulan el acceso a recursos compartidos, hacen observable la ejecución y comunican fallos temporales de forma clara.

El primer componente fue \textbf{JobQueue}, la cola de trabajos. Su regla de operación es simple de enunciar: existe un número máximo de inferencias que pueden ejecutarse a la vez y, cuando ese cupo está ocupado, cada nueva solicitud espera su turno en estricto orden de llegada. Se descartó un mecanismo de sincronización básico, como un semáforo, porque limita la concurrencia pero no garantiza el orden de llegada; en su lugar, la cola se construyó con las primitivas de sincronización de hilos de Python \cite{python_threading}, combinando una estructura de espera ordenada con señales que despiertan al trabajo más antiguo cuando se libera un cupo. Así se controla la concurrencia y se garantiza un trato equitativo entre solicitudes.

El límite de concurrencia se expone en la variable \textbf{MIDDLEWARE\_MAX\_CONCURRENT\_JOBS}, con valor predeterminado de 1. Funciona como un control indirecto de recursos: impide que varios backends carguen sus pesos en memoria al mismo tiempo y, en equipos con mayor capacidad, permite ampliar la concurrencia sin modificar el código.

El segundo componente fue \textbf{MetricsService}. Este servicio mantiene contadores protegidos contra accesos simultáneos: jobs recibidos, completados, fallidos, expirados por tiempo de espera y un conteo de errores por código. Los campos \textbf{active\_jobs} y \textbf{queued\_jobs} no son contadores acumulados, sino una fotografía del estado de la cola tomada en el momento de consultar \textbf{GET /api/v1/metrics}.

La integración con el servidor Uvicorn exigió un cuidado adicional: una solicitud de inferencia que espera su turno mantiene ocupado un hilo de atención, por lo que el servidor debe tolerar solicitudes en espera y, al mismo tiempo, seguir respondiendo los endpoints de gestión, como el de salud y el de métricas. Si se amplía el límite de concurrencia, la capacidad del servidor debe crecer en proporción \cite{uvicorn_docs}.

El tercer componente fue el estado \textbf{TIMEOUT}. Antes, toda excepción de inferencia terminaba en \textbf{FAILED}; sin embargo, un error funcional del backend y una inferencia que excede el tiempo máximo no representan el mismo problema. Con \textbf{TIMEOUT}, las métricas separan fallos funcionales de fallos por latencia, facilitando el diagnóstico sin revisar logs.

Cuando una inferencia supera el tiempo máximo permitido, el job termina en \textbf{TIMEOUT}, el contador de métricas correspondiente se incrementa y el middleware responde HTTP 504 con el código \textbf{DOCKER\_TIMEOUT}. Esta distinción se usó en las pruebas de carga y en el escenario de tiempo de espera forzado de la validación final.

\subsubsection{Fase 5 — Integración con la interfaz de ELAN}

La quinta fase trasladó las capacidades del middleware a la interfaz gráfica de ELAN. Hasta ese momento, la instalación y la inferencia podían ejecutarse desde Postman o clientes HTTP; sin embargo, el objetivo del TIC requería que el investigador permaneciera dentro del entorno de anotación. Por ello se implementaron acciones para instalar modelos, consultar conectividad, seleccionar el modelo activo, configurar la inferencia, ejecutar el reconocedor y mostrar errores desde ELAN \cite{auer2014avatech}.

La integración se ubicó en el paquete \textbf{mpi.\allowbreak{}eudico.\allowbreak{}client.\allowbreak{}annotator.\allowbreak{}recognizer.\allowbreak{}ai}. La clase \textbf{AIOrchestration\allowbreak{}MiddlewareClient} concentra la comunicación HTTP mediante \textbf{java.net.http.\allowbreak{}HttpClient} y fija \textbf{HTTP/1.1}, porque Uvicorn no ofrece HTTP/2 por defecto y durante las pruebas se identificaron fallos cuando el cliente intentaba negociar otra versión \cite{java11_httpclient,uvicorn_docs}. Esta clase evita duplicar solicitudes en el panel, el diálogo de modelos y el reconocedor.

El método \textbf{installModel} envía el ZIP en el formato \textbf{multipart/form-data}, el formato estándar con el que un cliente HTTP adjunta archivos \cite{rfc7578}. Como Java 11 no incluye un constructor nativo para ese tipo de cuerpo, el cliente lo arma manualmente respetando la estructura que el middleware espera al recibir archivos:

\begin{center}
\begin{minipage}{0.82\textwidth}
\hrule
\vspace{0.12cm}
\small
\begin{Verbatim}[baselinestretch=0.82]
--FormBoundary<uuid>
Content-Disposition: form-data; name="file"; filename="modelo.zip"
Content-Type: application/octet-stream

[bytes del archivo ZIP]
--FormBoundary<uuid>--
\end{Verbatim}
\vspace{0.05cm}
\hrule

\captionof{codigo}{Cuerpo multipart enviado por el cliente Java para instalar modelos}
\label{lst:multipart_install_model}
\end{minipage}
\end{center}

Un detalle de este formato provocó un error difícil de rastrear durante las pruebas: el delimitador declarado en la cabecera de la solicitud no debe llevar los dos guiones iniciales que sí llevan los delimitadores dentro del cuerpo. El diagnóstico de este problema y su solución quedaron registrados en la tabla de problemas frecuentes del Anexo~IV.

La clase \textbf{AIModel\allowbreak{}ManagerDialog} implementa el diálogo Swing de gestión de modelos. Incluye indicador de conexión, prueba de conectividad, tabla de modelos, instalación de ZIP, activación, desactivación y log de actividad. Las operaciones HTTP se ejecutan con \textbf{SwingWorker} para no bloquear el EDT de Swing; al finalizar, \textbf{done()} actualiza la interfaz en el hilo seguro para componentes gráficos \cite{java_swingworker}.

La clase \textbf{AIOrchestration\allowbreak{}RecognizerPanel} reemplazó parámetros hardcodeados por configuración visible: URL del middleware, modelo instalado, nivel de anotación destino, modo de etiquetado y video. Los modelos se cargan en segundo plano desde \textbf{addNotify()}, porque este método se ejecuta cuando el panel entra en la jerarquía visible y evita llamadas de red durante el constructor. Los valores sugeridos de nivel de anotación y modo se toman del campo \textbf{ui} del manifest.

La clase \textbf{AIOrchestration\allowbreak{}Recognizer} lee en tiempo de ejecución los valores seleccionados en el panel: modelo, nivel de anotación, modo de etiquetado y URL base. Con ello, la solicitud enviada al middleware refleja la configuración visible del usuario y no valores definidos en tiempo de compilación. También se crearon DTOs en \textbf{.ai.dto} para separar comunicación, representación de datos y lógica gráfica.

\textbf{SegmentVideoRequest} incorpora el método de fábrica \textbf{create}, que construye el payload con job, video, modelo, versión, nivel de anotación, etiqueta, modo de etiquetado, timeout de 300 segundos y \textbf{parameters} como objeto vacío. Esta decisión permite que el backend use su propia configuración interna y garantiza que el campo \textbf{parameters} siempre esté presente en la solicitud.

La comunicación entre \textbf{AIModel\allowbreak{}ManagerDialog} y \textbf{AIOrchestration\allowbreak{}RecognizerPanel} se resolvió con un flag booleano. Al cerrar el diálogo, el panel consulta \textbf{isModelListChanged()} y refresca el combo si hubo cambios. Finalmente, \textbf{addNotify()} dispara la carga inicial de modelos mediante \textbf{SwingWorker}; si el middleware no está disponible, el error aparece cuando el usuario visualiza el panel y puede interpretarlo en contexto.

\subsubsection{Fase 6 — Validación final}

La sexta fase verificó el sistema completo: middleware, Docker, backend del modelo y reconocedor integrado en ELAN. La validación combinó pruebas de caja negra y de caja blanca, porque no bastaba con probar endpoints aislados; había que asegurar la instalación, la activación, la selección en ELAN, la inferencia, el manejo de errores, las métricas y las anotaciones temporales. Las pruebas se realizaron con el modelo \textbf{lsec\_bio\_gloss\_final\_v1} y videos de LSEc, manteniendo la arquitectura real del sistema.

Se utilizó el modelo real en lugar de uno ficticio. Un modelo simulado habría validado la serialización y el transporte HTTP, pero no habría ejercitado la construcción de imágenes, el arranque de contenedores, la carga de pesos, la lectura de videos, la ejecución del pipeline ni la generación de segmentos reales. Esta decisión encareció las pruebas, aunque produjo evidencia representativa del escenario para el que se diseñó el middleware.

Los escenarios de pruebas se identifican con códigos E-XX. Estos códigos no corresponden a errores o estados HTTP, sino que son etiquetas internas de trazabilidad. Por ejemplo, E-07 corresponde a una inferencia exitosa desde ELAN; E-08, a la revisión visual de anotaciones en la línea de tiempo; y E-13, al tiempo de espera forzado. Esta nomenclatura relaciona la acción, el resultado esperado, la evidencia y la categoría de cada prueba.

Las pruebas de caja negra exploraron el comportamiento visible desde las interfaces expuestas al usuario final. Se definieron 15 escenarios en cuatro categorías: conectividad, gestión de modelos, inferencia y manejo de errores. Los escenarios cubrieron la comprobación de disponibilidad, la instalación, el listado y detalle de modelos, la activación y desactivación, la inferencia desde ELAN y desde Postman, la inserción de anotaciones en el nivel correcto y los fallos controlados: video inexistente, modelo desactivado, middleware detenido, contenedor detenido y tiempo de espera forzado, junto con la revisión de métricas. En cada caso se registraron la acción, la respuesta esperada y el comportamiento visible en ELAN, comprobando tanto la semántica de las respuestas HTTP como la experiencia del investigador.

Las pruebas de caja blanca revisaron estados internos que no se observan solo desde la API: la persistencia en \textbf{registry.json}, la generación del bootstrap manifest, la recuperación automática del registro, el incremento de métricas, la traducción de la ruta del video en los logs y la configuración del bind mount mediante \textbf{docker inspect}. Estas verificaciones confirmaron que el estado persistente, los montajes y los contadores internos quedaban consistentes después de cada operación.

Para rendimiento se usó \textbf{docker stats} en lugar de integrar \textbf{psutil}. La razón es que el consumo se distribuye entre dos contenedores: el middleware mantiene API, cola, contratos y orquestación; el backend ejecuta el pipeline, carga pesos y procesa video. Medir solo el proceso del middleware habría omitido el costo real de inferencia. \textbf{docker stats} permitió observar ambos contenedores desde el anfitrión y estimar el impacto sobre el hardware del usuario \cite{docker_cli_stats}. Los valores esperados para el equipo de desarrollo fueron menos de 5\% de CPU y menos de 150 MB de RAM para el middleware, y entre 60\% y 200\% de CPU junto con 1 a 4 GB de RAM para el backend.

El escenario \textbf{E-07} validó el flujo completo desde ELAN. El usuario selecciona el modelo, el video, el nivel de anotación destino y ejecuta el reconocedor. ELAN construye la solicitud, el middleware valida el contrato, localiza el modelo, traduce la ruta del video y solicita inferencia al backend. El resultado esperado es progreso visible y anotaciones en \textbf{AUTO\_GLOSS\_SEGMENTS}; sus tiempos deben coincidir con los intervalos visibles de señas. El escenario \textbf{E-08} complementa esta evidencia con la captura de la línea de tiempo donde se observan las anotaciones alineadas.

El escenario \textbf{E-13} validó timeout controlado. Desde Postman se envió una inferencia con \textbf{execution.timeout\_sec: 1}, valor menor al tiempo real de procesamiento. El resultado esperado fue HTTP 504 con \textbf{DOCKER\_TIMEOUT}. Después del error, el middleware debía seguir disponible, sin bloquear la cola, perder el registro ni dejar estados inconsistentes.

En la prueba de carga controlada se enviaron dos inferencias simultáneas con \textbf{MIDDLEWARE\_MAX\_CONCURRENT\_JOBS=1}. El comportamiento esperado era que la primera pasara a \textbf{RUNNING} y la segunda quedara en \textbf{QUEUED}; al terminar la primera, la segunda debía continuar en orden FIFO. Los logs registraron la transición de \textbf{QUEUED} a \textbf{RUNNING} y las métricas finales reportaron \textbf{total\_jobs: 2} con ambos trabajos completados. La prueba confirmó que la cola limita la concurrencia, preserva el orden de llegada y mantiene las métricas consistentes.


\subsection{Contratos de comunicación entre componentes}

El sistema se apoya en tres contratos de comunicación que ordenan el intercambio entre ELAN, el middleware y los backends de modelo. Estos contratos funcionan como una frontera clara entre tecnologías: ELAN mantiene su flujo de reconocedores AVATecH en Java, el middleware concentra la orquestación en Python y cada modelo ejecuta su inferencia dentro de un contenedor aislado. Por esta razón, los intercambios se definieron como cuerpos JSON sobre endpoints HTTP, una decisión coherente con el enfoque REST y con la necesidad de contar con interfaces verificables, documentables y desacopladas \cite{fielding2000rest}.

La estructura general de estos contratos se mantuvo estable durante el desarrollo. Los ajustes realizados en las iteraciones afectaron la validación, la compatibilidad o el manejo de errores, pero nunca la responsabilidad principal de cada componente. Esa estabilidad resultó clave para integrar ELAN, Docker y los backends de IA sin depender de los detalles internos de cada uno.

\subsubsection{Contrato entre ELAN y el Middleware (solicitud de inferencia)}

ELAN inicia el flujo mediante una solicitud \textbf{POST /api/v1/jobs/segment-video}. Esta operación no ejecuta reconocimiento dentro de ELAN; solo entrega al middleware los datos necesarios para identificar el trabajo, ubicar el video, seleccionar el modelo, definir el nivel de anotación destino y establecer un tiempo máximo de ejecución. Antes de iniciar operaciones costosas, el middleware valida la estructura recibida con modelos Pydantic integrados a FastAPI, de modo que errores de formato, campos incompletos o valores incompatibles se detecten en la entrada y no durante la ejecución del contenedor:

\begin{center}
\begin{minipage}{0.82\textwidth}
\hrule
\vspace{0.12cm}
\small
\begin{Verbatim}[baselinestretch=0.82]
{
  "job_id": "identificador-unico-del-job",
  "media": {"path": "C:/Users/.../video.mp4"},
  "model": {"model_id": "nombre_modelo", "version": "1.0.0"},
  "annotation": {
    "target_tier": "NombreTier",
    "default_label": "REGION",
    "label_mode": "gloss_top1"
  },
  "execution": {"timeout_sec": 300}
}
\end{Verbatim}
\vspace{0.05cm}
\hrule

\captionof{codigo}{Solicitud JSON enviada desde ELAN al middleware}
\label{lst:request_elan_middleware}
\end{minipage}
\end{center}

El campo \textbf{job\_id} permite seguir una solicitud en logs, métricas y consultas posteriores mediante \textbf{GET /api/v1/jobs/\{job\_id\}}. El objeto \textbf{media} contiene la ruta del video seleccionada desde ELAN, pero el middleware no la entrega al backend sin procesamiento, porque el contenedor ve una jerarquía de archivos distinta a la del anfitrión. El objeto \textbf{model} representa una versión concreta instalada, mientras que \textbf{annotation} define cómo se transformarán los segmentos en anotaciones: \textbf{target\_tier} indica el nivel de anotación destino, \textbf{default\_label} actúa como etiqueta de respaldo y \textbf{label\_mode} decide entre la glosa principal o una etiqueta simple. Finalmente, \textbf{execution.timeout\_sec} limita el tiempo de espera; si el backend lo supera, el middleware responde \textbf{DOCKER\_TIMEOUT} con HTTP 504.

\subsubsection{Contrato entre Middleware y el Backend de modelo (payload de inferencia)}

Después de validar la solicitud y seleccionar el modelo, el middleware traduce el contrato externo al contrato interno del backend. El cambio más importante ocurre en \textbf{media.path}: la ruta seleccionada en ELAN se convierte en la ruta visible dentro del contenedor según el bind mount de videos. Esta traducción es necesaria porque Docker aísla el sistema de archivos de cada contenedor y expone solo los directorios montados explícitamente \cite{docker2026overview}:

\begin{center}
\begin{minipage}{0.70\textwidth}
\hrule
\vspace{0.12cm}
\footnotesize
\begin{Verbatim}[baselinestretch=0.78]
{
  "job_id": "identificador-unico-del-job",
  "media": {"path": "/data/videos/video.mp4"},
  "annotation": { ... },
  "parameters": {}
}
\end{Verbatim}
\vspace{0.05cm}
\hrule

\captionof{codigo}{Payload interno enviado desde el middleware al backend}
\label{lst:payload_middleware_backend}
\end{minipage}
\end{center}

El campo \textbf{parameters} se incluye siempre, aunque sea un objeto vacío. Con ello, el backend puede leer una estructura uniforme y se evitan casos especiales por ausencia del campo. Al mismo tiempo, esto no significa que ELAN pueda sobrescribir libremente la configuración interna del modelo: el pipeline conserva su configuración base en \textbf{config/pipeline\_config.json}. Por tanto, \textbf{parameters} queda como un espacio de extensión controlado, sin reemplazar los artefactos ni las reglas declaradas en el paquete del modelo.

\subsubsection{Contrato entre el Backend de un modelo y el Middleware (respuesta de inferencia)}

Cuando termina la inferencia, el backend devuelve un JSON con la salida técnica del pipeline. El middleware no solo acepta una respuesta porque el código HTTP sea exitoso, sino que también verifica que el tipo de salida, el tiempo del video y una lista de segmentos con marcas de inicio y fin estén presentes. Esta validación evita que una respuesta incompleta llegue a ELAN y genere anotaciones inconsistentes en la línea de tiempo:

\begin{center}
\begin{minipage}{0.92\textwidth}
\hrule
\vspace{0.12cm}
\footnotesize
\begin{Verbatim}[baselinestretch=0.76]
{
  "output_type": "segments_with_gloss",
  "media_info": {"fps": 59.94, "duration_ms": 8508, "total_frames": 510},
  "segments": [{
    "start_ms": 884, "end_ms": 1518,
    "label": "TRABAJAR", "confidence": 0.164,
    "predictions": [
      {"rank": 1, "gloss_id": 10, "gloss": "TRABAJAR", "probability": 0.164},
      {"rank": 2, "gloss_id": 7, "gloss": "VER", "probability": 0.126}
    ]
  }]
}
\end{Verbatim}
\vspace{0.05cm}
\hrule

\captionof{codigo}{Respuesta JSON del backend con segmentos y glosas}
\label{lst:respuesta_backend_middleware}
\end{minipage}
\end{center}

El campo \textbf{output\_type} declara la forma semántica de la respuesta; en esta implementación, \textbf{segments\_with\_gloss} indica intervalos de tiempo acompañados de una etiqueta principal y de alternativas de predicción. El objeto \textbf{media\_info} aporta contexto del video procesado: \textbf{fps}, \textbf{duration\_ms} y \textbf{total\_frames}. Dentro de cada segmento, \textbf{label} es la etiqueta que se anotará, \textbf{confidence} refleja la confianza reportada por el backend y \textbf{predictions} lista las alternativas ordenadas por \textbf{rank}. La respuesta contiene así lo mínimo necesario para anotar en ELAN y, al mismo tiempo, información adicional útil para la revisión posterior. El Anexo~II presenta un ejemplo completo de esta respuesta, incluida la trazabilidad por etapas del campo \textbf{trace}.

La Tabla~\ref{tab:endpoints_api} resume los endpoints disponibles en la API REST del middleware, agrupados según su responsabilidad: disponibilidad del servicio, gestión de modelos, ejecución y consulta de inferencias, y lectura de métricas. Gracias a esta organización, ELAN consume únicamente los endpoints que su flujo necesita, mientras que herramientas como Postman o \textbf{curl} sirven de apoyo en la verificación durante las pruebas funcionales.

\begin{table}[H]
\centering
\caption{Endpoints de la API REST del middleware}
\label{tab:endpoints_api}
\small
\renewcommand{\arraystretch}{1.18}
\setlength{\tabcolsep}{5pt}

\begin{tabular}{
    >{\centering\arraybackslash}p{1.6cm}
    >{\raggedright\arraybackslash}p{5.0cm}
    >{\raggedright\arraybackslash}p{6.7cm}
}
\toprule
\textbf{Método} & \textbf{Endpoint} & \textbf{Descripción} \\
\midrule

\texttt{GET} &
\texttt{/health} &
Verifica la disponibilidad general del servicio. \\

\texttt{GET} &
\texttt{/api/v1/models} &
Lista todos los modelos instalados en el middleware. \\

\texttt{GET} &
\texttt{/api/v1/models/\{id\}} &
Recupera el detalle completo de un modelo específico. \\

\texttt{POST} &
\texttt{/api/v1/models/install} &
Instala un modelo a partir de un archivo ZIP enviado como \texttt{multipart/form-data}. \\

\texttt{PATCH} &
\texttt{/api/v1/models/\{id\}/status} &
Activa o desactiva un modelo previamente instalado. \\

\texttt{POST} &
\texttt{/api/v1/jobs/segment-video} &
Ejecuta una inferencia sobre un video seleccionado. \\

\texttt{GET} &
\texttt{/api/v1/jobs/\{job\_id\}} &
Recupera el resultado de un job de inferencia ejecutado. \\

\texttt{GET} &
\texttt{/api/v1/metrics} &
Consulta métricas acumuladas desde el inicio del servicio. \\

\bottomrule
\end{tabular}
\end{table}

\subsection{Despliegue y configuración del entorno}

El despliegue del middleware se realiza con Docker Compose porque el sistema necesita más que un proceso aislado: requiere una red compartida, comunicación con contenedores de modelos creados dinámicamente y persistencia entre reinicios. El archivo \textbf{docker-compose.yml} define el servicio principal, los bind mounts y la red común usada por los backends; su contenido completo, junto con las consideraciones de seguridad sobre el socket de Docker, consta en el Anexo~III. Esta decisión simplifica el arranque del sistema, ya que el usuario puede levantar el servicio con un comando único mientras Compose crea la red, publica el puerto y monta los directorios requeridos.

Se eligieron bind mounts para los datos persistentes porque dejan los archivos visibles en el sistema del anfitrión. Esto facilita auditar el registro de modelos, revisar manifests de respaldo y conservar información aunque se detengan o eliminen los contenedores. En cambio, los volúmenes con nombre pueden perderse con operaciones de limpieza de Docker, lo que no resulta conveniente para un entorno local de investigación \cite{docker2026overview}.

El middleware utiliza tres montajes con responsabilidades distintas. \textbf{./data/models\_store} conserva el registro y los archivos extraídos de los modelos instalados. \textbf{./data/bootstrap\_manifests} guarda copias de manifests para recuperar información mínima si el registro principal se pierde o queda inconsistente. Además, el socket de Docker permite que el SDK de Docker (para Python) construya imágenes, inicie contenedores, inspeccione estados y gestione backends desde el propio middleware \cite{docker_sdk_python}. Así, aunque el orquestador se ejecuta dentro de un contenedor, mantiene control local sobre los contenedores de inferencia.

La Tabla~\ref{tab:variables_entorno} describe las variables configurables del middleware. Estas variables separan la configuración del entorno respecto del código fuente, algo necesario porque cada estación de trabajo puede tener una ruta distinta para videos, diferente capacidad de hardware o distintos requisitos de logging.

\begin{table}[H]
\centering
\caption{Variables de entorno del middleware}
\label{tab:variables_entorno}
\scriptsize
\renewcommand{\arraystretch}{1.18}
\setlength{\tabcolsep}{5pt}

\begin{tabular}{
    >{\raggedright\arraybackslash}p{4.2cm}
    >{\raggedright\arraybackslash}p{3.4cm}
    >{\raggedright\arraybackslash}p{5.8cm}
}
\toprule
\textbf{Variable} & \textbf{Valor por defecto} & \textbf{Descripción} \\
\midrule

\texttt{MIDDLEWARE\_HOST} &
\texttt{0.0.0.0} &
Interfaz de escucha dentro del contenedor. \\

\texttt{MIDDLEWARE\_PORT} &
\texttt{8000} &
Puerto expuesto al host. \\

\texttt{MIDDLEWARE\_LOG\_LEVEL} &
\texttt{INFO} &
Nivel de logging estructurado. \\

\texttt{MIDDLEWARE\_MODELS\_}\newline
\texttt{STORE\_DIR} &
\texttt{/app/data/}\newline
\texttt{models\_store} &
Directorio donde se almacena el registry de modelos. \\

\texttt{MIDDLEWARE\_BOOTSTRAP\_}\newline
\texttt{MANIFESTS\_DIR} &
\texttt{/app/data/bootstrap/}\newline
\texttt{manifests} &
Directorio persistente para los manifests de respaldo. \\

\texttt{MIDDLEWARE\_VIDEOS\_DIR} &
\textit{requerido} &
Ruta del host donde se encuentran los videos disponibles para inferencia. \\

\texttt{MIDDLEWARE\_DOCKER\_}\newline
\texttt{NETWORK} &
\texttt{elan-ai-shared} &
Red Docker compartida entre el middleware y los backends de los modelos. \\

\texttt{MIDDLEWARE\_MAX\_}\newline
\texttt{CONCURRENT\_JOBS} &
\texttt{1} &
Límite de inferencias simultáneas permitidas por el middleware. \\

\bottomrule
\end{tabular}
\end{table}

La variable \textbf{MIDDLEWARE\_VIDEOS\_DIR} es la única requerida sin valor por defecto porque depende del equipo donde se despliega el sistema. Debe apuntar al directorio del anfitrión que contiene los videos de LSEc disponibles para inferencia. El middleware no intenta inferir esta ruta: la entrega al SDK de Docker como fuente del bind mount usado por los contenedores de modelo. Esta decisión evita transformaciones ambiguas entre sistemas operativos y se relaciona con el problema resuelto en la Fase 3, donde ELAN trabaja con rutas del anfitrión y el backend con rutas internas del contenedor.

Conviene señalar que el despliegue local de un solo usuario es el escenario previsto del TIC, pero no un límite de la arquitectura. Por diseño, el puerto del middleware se publica únicamente en la interfaz de loopback (\textbf{127.0.0.1}), de modo que solo el ELAN de la misma máquina puede consumirlo. Sin embargo, al tratarse de una relación cliente-servidor sobre HTTP, bastaría con publicar el puerto en la dirección IP del equipo dentro de una red local para que varias estaciones de trabajo con ELAN enviaran solicitudes al mismo middleware, configurando esa dirección en el campo de URL del panel del reconocedor. En ese caso, la cola FIFO atendería las solicitudes en orden de llegada y la variable \textbf{MIDDLEWARE\_MAX\_CONCURRENT\_JOBS} determinaría cuántas inferencias se ejecutan en paralelo según la capacidad del servidor; con el valor por defecto de 1, las peticiones de los demás clientes esperarían su turno. Este escenario, junto con la posibilidad de trasladar el middleware y los modelos a un servidor de alto rendimiento o a la nube, se retoma en las recomendaciones del capítulo de resultados.

Para levantar el sistema por primera vez, o después de cambios en el código, se ejecutan los comandos mostrados a continuación. El primer comando construye la imagen del middleware y deja el servicio en segundo plano; el segundo verifica que el endpoint de salud responda antes de instalar modelos o ejecutar inferencias:

\begin{center}
\begin{minipage}{0.78\textwidth}
\hrule
\vspace{0.10cm}
\footnotesize
\begin{Verbatim}[baselinestretch=0.78]
# Desde el directorio middleware/
docker compose up --build -d

# Verificar que el servicio está activo
curl http://localhost:8000/health
\end{Verbatim}
\vspace{0.03cm}
\hrule

\captionof{codigo}{Comandos iniciales para desplegar y verificar el middleware}
\label{lst:comandos_despliegue_middleware}
\end{minipage}
\end{center}

La integración con ELAN requiere compilar el árbol fuente modificado con Maven e incorporar el JAR resultante a la instalación de la herramienta; el procedimiento paso a paso, incluido el registro del reconocedor mediante el mecanismo SPI de Java, se describe en el Anexo~IV. Esta integración no reemplaza el flujo de anotación existente; añade un reconocedor dentro del mecanismo de extensión previsto por AVATecH \cite{auer2014avatech}. Una vez instalado, el reconocedor \textit{AI Orchestration Middleware} aparece en la lista de reconocedores de ELAN y permite seleccionar modelo, configurar el nivel de anotación destino y ejecutar inferencias sin manipular directamente los contenedores Docker.

\vspace{1\baselineskip}

\subsection{Estructura de directorios del sistema}

El Código~\ref{lst:estructura_directorios} muestra cómo se organizó el middleware para reflejar la separación de responsabilidades definida en el diseño. La capa \textbf{api/} recibe solicitudes HTTP y serializa respuestas; \textbf{core/} centraliza configuración y logging; \textbf{schemas/} contiene los modelos Pydantic que hacen ejecutables los contratos; \textbf{services/} agrupa la lógica de instalación, registro, jobs, Docker, cola y métricas; \textbf{runners/} abstrae el mecanismo de ejecución; \textbf{storage/} conserva jobs en memoria durante la sesión; y \textbf{data/} mantiene información persistente mediante bind mounts.

\begin{center}
\begin{minipage}{0.96\textwidth}
\hrule
\vspace{0.10cm}
\scriptsize
\begin{Verbatim}[baselinestretch=0.72]
middleware/
|-- app/
|   |-- main.py                    (FastAPI: inicio, routers, error handlers)
|   |-- api/
|   |   |-- routes_health.py       (GET /health)
|   |   |-- routes_models.py       (CRUD modelos e instalación)
|   |   |-- routes_jobs.py         (POST /jobs/segment-video)
|   |   `-- routes_metrics.py      (GET /metrics)
|   |-- core/
|   |   |-- config.py              (Variables de entorno -> Settings)
|   |   `-- logging_config.py      (Logging estructurado)
|   |-- runners/
|   |   |-- base_runner.py         (Clase abstracta BaseRunner)
|   |   |-- docker_runner.py       (DockerRunner: runtime activo)
|   |   `-- runner_selector.py     (Selección de runner por manifest)
|   |-- schemas/                   (Contratos Pydantic)
|   |-- services/
|   |   |-- model_registry_service.py  (Instalar y gestionar modelos)
|   |   |-- job_service.py             (Orquestar inferencia completa)
|   |   |-- docker_service.py          (Wrapper del SDK Docker)
|   |   |-- docker_lifecycle_service.py (Build, start y comprobación de disponibilidad)
|   |   |-- job_queue.py               (Cola FIFO y concurrencia)
|   |   `-- metrics_service.py         (Contadores thread-safe)
|   `-- storage/
|       `-- memory_store.py        (Jobs en memoria por sesión)
|-- data/
|   |-- models_store/              (Registry y archivos de modelos)
|   `-- bootstrap_manifests/       (Manifests de respaldo)
|-- Dockerfile                     (Imagen del middleware)
|-- docker-compose.yml             (Orquestación del despliegue)
`-- requirements.txt               (Dependencias Python)
\end{Verbatim}
\vspace{0.03cm}
\hrule

\captionof{codigo}{Estructura de directorios del middleware}
\label{lst:estructura_directorios}
\end{minipage}
\end{center}

Los archivos Java añadidos a ELAN se ubicaron en el paquete \textbf{mpi.eudico.client.annotator.recognizer.ai}, como subpaquete del módulo de reconocedores. Esta ubicación mantiene el nuevo código cerca del punto de extensión AVATecH y, al mismo tiempo, separa la integración con IA de los reconocedores previos. El Código~\ref{lst:estructura_clases_elan} muestra esta organización, donde el cliente HTTP, el reconocedor, el panel, el diálogo de modelos y los DTO principales quedan agrupados de forma explícita; el inventario completo de clases y DTO consta en el Anexo~IV.

\begin{center}
\begin{minipage}{0.96\textwidth}
\hrule
\vspace{0.10cm}
\scriptsize
\begin{Verbatim}[baselinestretch=0.72]
src/main/java/mpi/eudico/client/annotator/recognizer/ai/
|-- AIOrchestrationMiddlewareClient.java    (cliente HTTP del middleware)
|-- AIOrchestrationMiddlewareException.java (excepción con HTTP status y body)
|-- AIOrchestrationRecognizer.java          (reconocedor AVATecH)
|-- AIOrchestrationRecognizerPanel.java     (panel de configuración en ELAN)
|-- AIModelManagerDialog.java               (diálogo de gestión de modelos)
`-- dto/
    |-- ModelSummary.java                   (resumen de modelo para lista)
    |-- InstalledModelDetail.java           (detalle completo con UI config)
    |-- ModelUiConfig.java                  (configuración UI del modelo)
    |-- SegmentVideoRequest.java            (solicitud de inferencia)
    `-- SegmentVideoResponse.java           (respuesta con segmentos)
\end{Verbatim}
\vspace{0.03cm}
\hrule

\captionof{codigo}{Estructura de clases Java integradas en ELAN}
\label{lst:estructura_clases_elan}
\end{minipage}
\end{center}

\vspace{1\baselineskip}

\subsection{Guía de replicación: creación de un nuevo backend de modelo}

Una propiedad central del sistema es que admite nuevos modelos sin modificar el código de ELAN ni del middleware. Para lograrlo, cada modelo empaqueta su propia variabilidad: dependencias, artefactos, configuración, forma de ejecución y contrato HTTP. Si el paquete ZIP cumple el contrato descrito en la Fase 2 y detallado en el Anexo~I, el middleware puede instalarlo, construir su imagen, iniciar el contenedor y mostrarlo en ELAN con el mismo flujo usado por los modelos ya validados.

El primer paso es implementar el backend como un servicio HTTP que exponga los endpoints requeridos. Puede usarse FastAPI u otro framework, porque el middleware no depende de la implementación interna del modelo; solo necesita que el contenedor cumpla el contrato. \textbf{GET /health} debe responder HTTP 200 únicamente cuando el pipeline esté listo para inferir, ya que un contenedor iniciado no siempre significa que pesos, vocabularios o recursos de visión artificial estén cargados. \textbf{POST /infer} debe aceptar el payload definido en la sección de contratos y responder con \textbf{output\_type}, \textbf{media\_info} y \textbf{segments}. Cada segmento debe incluir al menos \textbf{start\_ms}, \textbf{end\_ms} y \textbf{label}; los campos \textbf{confidence} y \textbf{predictions} pueden añadirse cuando el modelo entregue evidencia adicional.

El segundo paso es empaquetar el backend con sus pesos y configuración, colocando \textbf{manifest.json} en la raíz del paquete. Este manifest identifica el modelo, su versión, los artefactos requeridos y la forma de construir y ejecutar el backend. El \textbf{Dockerfile} debe instalar las dependencias necesarias y definir cómo inicia el servicio HTTP. Esta encapsulación permite aislar bibliotecas de IA potencialmente incompatibles con ELAN o con otros modelos, manteniendo el principio de contenedorización de la arquitectura \cite{docker2026overview}. Los artefactos, como pesos, vocabularios y archivos de configuración, se declaran con rutas relativas dentro del ZIP para que el paquete pueda instalarse en distintos equipos.

El tercer paso es comprimir la estructura de directorios en un ZIP, cuidando que \textbf{manifest.json} quede en la raíz y no dentro de un subdirectorio adicional. Esta condición permite que el instalador valide el contrato antes de extraer y construir el backend. En Windows PowerShell, el comando correcto es:

\begin{center}
\begin{minipage}{0.78\textwidth}
\hrule
\vspace{0.08cm}
\footnotesize
\begin{Verbatim}[baselinestretch=0.78]
cd carpeta_que_contiene_el_paquete\
Compress-Archive -Path nombre_paquete\* -DestinationPath modelo.zip
\end{Verbatim}
\vspace{0.03cm}
\hrule

\captionof{codigo}{Comando PowerShell para comprimir un paquete de modelo}
\label{lst:comprimir_modelo_zip}
\end{minipage}
\end{center}

El cuarto paso es instalar el paquete mediante \textbf{POST /api/v1/models/install} o desde ELAN con el botón \textit{Instalar ZIP}. En ambos casos, el flujo técnico es el mismo: el middleware recibe el ZIP como \textbf{multipart/form-data}, localiza y valida el manifest, verifica artefactos, extrae el paquete, actualiza el registro, construye la imagen Docker y ejecuta las comprobaciones de arranque. Si todo es correcto, el modelo queda instalado, el contenedor se inicia y aparece disponible en el panel del reconocedor.

Este proceso concentra el esfuerzo de integración en el backend y en la preparación correcta del paquete. El desarrollador no necesita conocer el código interno de ELAN ni del middleware; debe respetar el contrato de \textbf{/health}, \textbf{/infer} y \textbf{manifest.json}. Con ello se conserva el desacoplamiento de la arquitectura: ELAN no conoce las bibliotecas del modelo, el middleware no interpreta la lógica interna del pipeline y el backend no manipula archivos EAF ni estructuras propias de ELAN.

\vspace{1\baselineskip}

\subsection{Uso ético de herramientas de inteligencia artificial generativa}

El presente Trabajo de Integración Curricular es de autoría propia. Durante su desarrollo se utilizaron de manera ética y responsable herramientas de inteligencia artificial generativa (ChatGPT y Claude) como apoyo técnico, en concordancia con las orientaciones internacionales sobre el uso de esta tecnología en contextos de investigación \cite{unesco2023guidance}.

Su aporte principal se dio en la etapa de ingeniería inversa del código fuente de ELAN 7.1. Dado que la herramienta carece de documentación interna detallada sobre su mecanismo de reconocedores, estas herramientas se emplearon para analizar el paquete \textbf{mpi.eudico.client.annotator.recognizer}, ubicar los puntos de extensión del protocolo AVATecH y orientar la forma de implementar las interfaces requeridas, como \textbf{Recognizer2}, \textbf{IncrementalRecognizer} y la comunicación mediante \textbf{RecognizerHost}. También sirvieron de apoyo puntual para interpretar errores de integración durante las pruebas y para revisar fragmentos repetitivos de código del prototipo.

Si bien la asistencia de IA facilitó la comprensión inicial de un código fuente extenso y ajeno, cada sugerencia debió verificarse contra el código real de ELAN 7.1 y la especificación de AVATecH, ajustarse al comportamiento observado durante las pruebas de integración y, en varios casos, descartarse por no corresponder a la versión utilizada de la herramienta. Este ejercicio de contraste fortaleció mis habilidades de lectura de código de terceros, diseño de integraciones entre tecnologías heterogéneas y análisis crítico de los resultados generados por estas herramientas.

Las decisiones metodológicas, los ajustes técnicos, la validación de resultados y las conclusiones corresponden exclusivamente a mi autoría.
\vspace{1\baselineskip}